
\subsection{Accrue current tax before choosing its payment}

A conditional calculation supplies an amount, but recognition and payment
change different financial balances. The next finite profile makes those
changes explicit. It begins with the complete recorded return-and-lease
calculation, recomputes the conditional loss result and requires an affirmative
recognition premise. An unknown premise remains unresolved; a false premise
does not produce a financial state. The profile applies only when the native
period has no current-tax expense already recorded and its tax-close cursor
still equals the assessment start. It therefore does not guess how an opening
tax accrual or an earlier assessment should be allocated.

The successor adds a separate income-tax payable to the original 25 accounts.
That liability is not trade payable. If the recomputed gross conditional tax
after loss usage is $C$, the accrual debits current-tax expense and credits
income-tax payable by $C$. Cash does not change. A zero amount produces a zero
journal. Existing deferred-tax assets, deferred-tax liabilities and deferred-tax
expense remain exactly as recorded; the loss balance does not create a new
deferred asset.

Payment requires its own identity, date, positive amount and affirmative
instruction. It debits income-tax payable and credits cash, subject to both the
outstanding liability and available cash. The author supplies the amount;
neither full payment nor a refund is inferred. A shortage of cash blocks that
payment while leaving a valid accrual available. This illustrative profile
classifies its specified tax payments as operating cash movements.

\begin{samepage}
Let $S$ be cumulative payment against the newly accrued amount $C$, with
$0\leq S\leq C$. Relative to the original recorded statements, the changes are
\[
 \Delta\text{profit}=-C,\qquad
 \Delta\text{tax payable}=C-S,\qquad
 \Delta\text{cash}=-S.
\]
\end{samepage}
Assets decrease by $S$, liabilities increase by $C-S$ and equity decreases by
$C$, preserving the balance-sheet identity. In the indirect operating cash
bridge, the lower profit and the increase in tax payable contribute
$-C+(C-S)=-S$. This equals the separately recorded cash payment. Paying the
liability does not charge the same expense again.

For the separate hypothetical 2026 calculation, loss usage leaves conditional
tax of $1/2$. Accrual records expense and payable of $1/2$ while cash stays at
500. An explicit payment of $1/4$ leaves cash of $1999/4$ and payable of $1/4$;
a second payment of $1/4$ clears the payable and leaves cash of $999/2$.
The recorded deferred-tax asset of 10 and liability of 4 are unchanged.
The original mixed example instead generates the loss quantity 19 and a zero
tax journal. These are consequences of the declared synthetic inputs, not
assessments of either a company or a jurisdiction.

The financial date and assessment coverage have separate meanings. Accrual is
dated on the last included day; coverage ends at the following exclusive
endpoint. The successor records that endpoint separately and leaves the native
tax-close cursor unchanged. Payments cannot precede the last tax operation or
fall outside the original financial period. An end-of-period accrual therefore
cannot silently accept a payment in the next year: financial-period carry must
first be specified.

The original native schema refuses the additional account and the exclusive
close cursor, so this is a distinct financial profile. Its validation replays
the retained calculation and every tax payment, then compares the complete
state. Changed balances, dates, loss results, payment histories and proposal
flags are rejected. Repeating the same accrual request produces the same pure
proposal; it does not establish globally unique durable posting. Later ordinary
return or lease events, financial-year carry, actual accounting admission and
annual completeness remain separate requirements.
